% ===========================================================================
\part{Architecture}
% ===========================================================================

\chapter{Layers, and the shape to grow into}
\label{ch:layers}

\section{The organising principle: vertical slices, horizontal gates}

You already made a good structural decision without necessarily framing it as one. The code is
organised by \emph{subject} --- \file{portfolio/}, \file{publications/},
\file{job\_experience/}, \file{personal\_information/} --- rather than by \emph{technical kind}
(a \file{models/} directory, a \file{repositories/} directory). That is the right instinct. Each
subject is a vertical slice, and everything about publications lives in one place, so adding a
field to a publication touches one directory.

What is missing is the second axis. Within each slice, code needs to be separated by \emph{which
of the two programs it belongs to}. Right now that separation happens at the top of
\file{src/lib.rs}, at whole-module granularity, which is too coarse: it forces an entire subject to
be server-only when only half of it needs to be.

\begin{keyidea}
Organise \emph{directories} by subject. Organise \emph{files within a directory} by which side of
the boundary they live on. The \code{cfg} gates then sit on individual files in
\file{<subject>/mod.rs}, not on whole subjects in \file{lib.rs}.
\end{keyidea}

\section{The target tree}

Here is the layout to grow into. Compare it against the inventory in Chapter~\ref{ch:state}; most
of it is a redistribution of files you already have.

\begin{lstlisting}[style=tree,caption={The target module tree. \textcolor{rust}{\emph{ssr}} marks
files compiled only into the server; everything else is compiled into both programs.}]
src/
  lib.rs                      # module roots. Gates live here only for genuinely server-only subjects.
  main.rs                     # the composition root: config, pool, Actix. ssr only by construction.
  app.rs                      # <App/>: the shell, <Routes>, global layout
  error.rs                    # AppError: the one error type the browser is allowed to see

  ui/                         # subject-AGNOSTIC presentation only
    mod.rs
    layout.rs                 #   header, nav, footer
    components.rs             #   reusable pieces: Card, TagList, DateRange, ExternalLink
    pages/                    #   one file per route; a page composes several subjects

  portfolio/                  # a vertical slice
    mod.rs                    #   declares the files below, and holds the cfg gates
    dto.rs                    #   PortfolioItemDto, TagDto        -- crosses the wire
    api.rs                    #   #[server] list_portfolio(), get_portfolio_item()
    components.rs             #   <PortfolioList/>, <PortfolioCard/>  -- subject-specific
    model.rs                  # ssr  Diesel structs: Queryable, Selectable
    repository.rs             # ssr  the SQL
  publications/               # same six files
  job_experience/             # same six files
  personal_information/       # same six files

  schema.rs                   # ssr  diesel table! macros, generated -- never edit by hand
  database/
    mod.rs
    connection.rs             # ssr  build the r2d2 pool from a config value (not from env)
  config.rs                   # ssr  typed AppConfig, loaded once in main
\end{lstlisting}

Three things changed relative to today, and each one has a reason.

\begin{description}
  \item[\file{src/entrypoint/} narrowed to composition, not controllers.] The original
    \file{routes/request/response} shape was reaching for a controller layer that this stack
    generates for you: server functions \emph{are} the inbound adapter, and \file{api.rs} inside
    each slice is where they live. Writing Actix handlers that return JSON for your own front end
    duplicates \code{\#[server]} and loses type safety across the wire.

    That subtree has since been deleted in favour of \file{http.rs} and
    \file{application\_state.rs}, which is a much better use of the directory --- see
    \S\ref{sec:appstate}. The remaining risk is only that two declared-but-empty modules are an
    invitation; fill them or remove them, rather than leaving them indefinitely. Defect \dnum{2}.

  \item[Each slice gained \file{dto.rs} and \file{api.rs}, and both are ungated.] This is the
    single most important structural change in the document. Chapter~\ref{ch:dto} explains the
    two-struct rule and Chapter~\ref{ch:composition} shows the wiring; the short version is that
    a \code{\#[server]} function \emph{must} be visible to the browser build, so the module
    holding it cannot be gated.

  \item[\file{src/ui/} is new.] Components do not belong in \file{app.rs}. \file{app.rs} should
    stay small enough to read in one screen: it is a table of contents for the site, not a place
    where markup accumulates.
\end{description}

\begin{pitfall}[The gate that will bite you]
\code{\#[cfg(feature = "ssr")] pub mod publications;} makes every server function inside
\file{publications/} invisible to the browser build. The error is
\code{error[E0433]: failed to resolve: could not find `publications` in the crate root}, and
rustc helpfully adds \code{found an item that was configured out ... the item is gated behind the
`ssr` feature}. The fix is never to un-gate the \emph{query} code --- it is to move the gate down
one level, onto the files that actually contain Diesel. This is defect \dnum{5}.
\end{pitfall}

\begin{verified}[This is happening right now, in \file{src/publications/}]
While this document was being written, a component appeared in the repository:

\begin{lstlisting}[style=rs]
// src/publications/mod.rs
pub mod components;      // <-- new
pub mod model;
pub mod repository;

// src/publications/components.rs
use leptos::prelude::*;

#[component]
fn Publications() -> impl IntoView {
    view! { <p> "Aca van las publicaciones" </p> }
}
\end{lstlisting}

This is a good instinct --- keeping the component next to the data it renders --- and it has walked
straight into the gate. Two things are wrong, and neither produces an error yet:

\textbf{The component is in the server-only build.} \file{src/lib.rs} still says
\code{\#[cfg(feature = "ssr")] pub mod publications;}, so \code{components.rs} is compiled into
the server and \emph{not} into the wasm bundle. It compiles today only because nothing references
it. The moment \file{app.rs} adds \code{<Publications/>}, the browser build fails with the
\code{E0433} above. Moving the gate off \code{publications} and onto \code{model} and
\code{repository} fixes it --- and that gate move is exactly the change in
\file{<subject>/mod.rs} shown below.

\textbf{The component is private.} \code{fn Publications()} has no \code{pub}, so even in the
server build nothing outside \file{components.rs} can name it. Components that a route renders
need to be \code{pub}.
\end{verified}

\begin{why}[Components in the slice, or in \file{ui/}?]
Both are defensible, and the tree above shows only one of them. Co-locating a component with its
data --- \file{publications/components.rs} --- keeps a feature in one directory and is easier to
delete cleanly. A shared \file{ui/} directory keeps the pieces that are genuinely reusable
(\code{Card}, \code{TagList}, \code{DateRange}) from being buried inside whichever subject
happened to need them first.

The split that works in practice: \textbf{subject-specific components live in the subject}
(\file{publications/components.rs} renders publications and nothing else);
\textbf{subject-agnostic components live in \file{ui/}}. Pages --- the things a route maps to ---
go in \file{ui/pages/}, because a page usually composes more than one subject. The one hard rule is
the gate: any file containing a component or a \code{\#[server]} function must be reachable in
both builds.
\end{why}

\subsection{What \file{<subject>/mod.rs} looks like}

\begin{lstlisting}[style=rs,caption={\file{src/publications/mod.rs} --- the gate lives here, per file. This
is the change that unblocks the component already sitting in that directory.}]
pub mod api;                                  // #[server] functions: BOTH programs
pub mod components;                           // <Publications/>:     BOTH programs
pub mod dto;                                  // wire types:          BOTH programs

#[cfg(feature = "ssr")]
pub mod model;                                // Diesel structs: server only
#[cfg(feature = "ssr")]
pub mod repository;                           // SQL:           server only
\end{lstlisting}

and correspondingly \file{src/lib.rs} loses four of its seven gates:

\begin{lstlisting}[style=rs,caption={\file{src/lib.rs} --- only the genuinely server-only roots stay gated.}]
pub mod app;
pub mod error;
pub mod ui;
pub mod job_experience;                       // ungated: each holds DTOs, server fns, components
pub mod personal_information;
pub mod portfolio;
pub mod publications;

#[cfg(feature = "ssr")] pub mod config;
#[cfg(feature = "ssr")] pub mod database;
#[cfg(feature = "ssr")] pub mod schema;
\end{lstlisting}

An equivalent and equally idiomatic convention --- the one the upstream Leptos examples use --- is
to keep a slice in a single file with an inner gated module:

\begin{lstlisting}[style=rs,caption={The single-file variant, for slices small enough not to need four files.}]
// src/publications/mod.rs
use crate::error::AppError;
use serde::{Deserialize, Serialize};

#[derive(Debug, Clone, Serialize, Deserialize)]
pub struct PublicationDto { /* ... */ }

#[cfg(feature = "ssr")]
mod ssr {                                     // everything server-only, in one place
    pub use crate::database::connection::DbPool;
    pub use crate::schema::publication_items;
    pub use diesel::prelude::*;
}

#[server]
pub async fn list_publications() -> Result<Vec<PublicationDto>, AppError> {
    #[cfg(feature = "ssr")]
    use self::ssr::*;                         // the body is compiled only for ssr anyway
    // ...
}
\end{lstlisting}

Pick one convention and hold to it. Four files per slice is better once a slice has more than one
query and more than one DTO, which all four of yours will.

\section{Dependency direction}

Layers are only useful if the arrows point one way. Figure~\ref{fig:layers} states the rule for
this project.

\begin{diagram}{Permitted dependency directions. Every arrow points \emph{down or inward}. The two
crossings marked in red are the ones to watch for --- both are easy to write and both invert the
architecture.}{fig:layers}
\begin{tikzpicture}[node distance=5mm]
  \node[cli,minimum width=64mm] (ui) {\textbf{\file{ui/}} --- pages and components\\[1pt]\scriptsize
    knows DTOs and server-fn names. Knows nothing else.};
  \node[shared,minimum width=64mm,below=6mm of ui] (api) {\textbf{\file{<subject>/api.rs}} --- server functions\\[1pt]\scriptsize
    the inbound adapter: extract state, call the repository, map to DTO};
  \node[shared,minimum width=64mm,below=6mm of api] (dto) {\textbf{\file{<subject>/dto.rs}} $+$ \file{error.rs}\\[1pt]\scriptsize
    plain serde types. Depend on nothing in this project.};
  \node[srv,minimum width=64mm,below=9mm of dto] (rep) {\textbf{\file{<subject>/repository.rs}}\\[1pt]\scriptsize
    owns the pool and the SQL};
  \node[srv,minimum width=64mm,below=6mm of rep] (mod) {\textbf{\file{<subject>/model.rs}} $+$ \file{schema.rs}\\[1pt]\scriptsize
    Diesel's view of the tables};

  \draw[flow] (ui) -- (api);
  \draw[flow] (api) -- (dto);
  \draw[flow] (api.east) -- ++(9mm,0) |- (rep.east);
  \draw[flow] (rep) -- (mod);
  \draw[flow] (rep.west) -- ++(-9mm,0) |- (dto.west);

  % forbidden
  \draw[-{Stealth[length=4pt]},line width=0.9pt,draw=crimson,dashed]
    (mod.west) -- ++(-16mm,0) |- node[lbl,pos=0.75,left,align=right,text=crimson]{a Diesel model\\in a \code{view!}}(ui.west);
  \draw[-{Stealth[length=4pt]},line width=0.9pt,draw=crimson,dashed]
    (dto.east) -- ++(20mm,0) |- node[lbl,pos=0.72,right,align=left,text=crimson]{a DTO that\\imports \code{schema}}(mod.east);

  \begin{scope}[on background layer]
    \node[lane,fit=(ui)(dto),inner sep=3mm,label={[tag,anchor=north west]north west:compiled into both programs}] {};
    \node[lane,fit=(rep)(mod),inner sep=3mm,label={[tag,anchor=south west]south west:server only}] {};
  \end{scope}
\end{tikzpicture}
\end{diagram}

The two red arrows are worth naming, because they are the two ways this design decays:

\begin{enumerate}
  \item \textbf{A Diesel model rendered directly in a \code{view!}.} It compiles, if you also
    un-gate the model, and then \code{diesel} follows it into the wasm build and the build breaks
    with errors about \code{mio} or missing \code{core} items. That failure mode is documented in
    the Leptos book precisely because it is common; \code{diesel} with the \code{postgres} feature
    links native \code{libpq}, which has no WebAssembly target.
  \item \textbf{A DTO that reaches for \file{schema.rs}} --- usually because someone added
    \code{\#[diesel(table\_name = ...)]} to the DTO to avoid writing a conversion. That drags the
    schema into the shared layer and the boundary is gone.
\end{enumerate}

\section{One crate or a workspace?}

Worth deciding now, because migrating later is tedious. Both shapes are supported by
cargo-leptos.

\begin{table}[htbp]
\centering
\caption{The two project shapes, judged for this project specifically.}
\label{tab:shapes}
\small
\begin{tabular}{@{}L{0.14\textwidth}L{0.40\textwidth}L{0.40\textwidth}@{}}
\toprule
 & \textbf{Single crate (today)} & \textbf{Workspace: \code{shared} / \code{app} / \code{frontend} / \code{server}} \\
\midrule
Boundary enforced by & discipline plus \code{cfg} gates & the compiler --- \code{server} is never built for wasm \\
Config & \code{[package.metadata.leptos]} & \code{[[workspace.metadata.leptos]]} with \code{bin-package} and \code{lib-package} \\
Cost & one \code{Cargo.toml}; every gate is yours to get right & four manifests, and the shared crate \emph{still} needs \code{ssr}/\code{hydrate} gates \\
Risk & a mis-gated module reaches wasm & \code{cargo build --workspace} unifies features and silently enables \code{ssr}+\code{hydrate} together \\
\bottomrule
\end{tabular}
\end{table}

\begin{pattern}[Stay with one crate]
For a personal website with four content types, the single-crate layout is the right call. The
workspace's benefit is that it makes the boundary a compiler-enforced fact rather than a
convention --- genuinely valuable at scale, and not worth four manifests here. The gates in
\file{lib.rs} plus the two-build check from Chapter~\ref{ch:twobuilds} give you most of the safety
for a fraction of the ceremony.

Revisit this if either becomes true: the server grows logic you want to unit-test without
compiling wasm at all, or the wasm bundle grows past roughly a megabyte and you need to audit what
is in it.
\end{pattern}

\begin{pitfall}[If you ever do move to a workspace]
Two traps, both silent. \code{[[workspace.metadata.leptos]]} is an \emph{array} of tables ---
single brackets and cargo-leptos simply will not find the project. And profile sections
(\code{[profile.wasm-release]}) are honoured only in the \emph{root} manifest; put them in a
member and your \code{opt-level = 'z'} is ignored with a warning you will not read.
\end{pitfall}

% ===========================================================================
\chapter{Domain models and DTOs: the two-struct rule}
\label{ch:dto}

\section{The rule}

\begin{keyidea}
For every content type, write two structs. One describes the \emph{table} and lives on the server.
One describes the \emph{payload} and travels to the browser. Convert between them with
\code{From}. Do not try to make one struct do both jobs.
\end{keyidea}

This will feel like duplication for about a week, and then it will start paying for itself. The
reasons are not aesthetic.

\begin{why}
\textbf{They have different lifetimes.} The table struct changes when you run a migration; the
payload changes when the page needs different information. Coupling them means every schema tweak
is a wire-format change.

\textbf{They have different derives, and the derives are mutually hostile.} The table struct needs
\code{Queryable} and \code{Selectable} from Diesel, which do not exist in the browser build. The
payload needs \code{Serialize} and \code{Deserialize} present in \emph{both} builds. One struct
cannot satisfy both sets without \code{cfg\_attr} gymnastics that break in a way described below.

\textbf{They have different obligations about truth.} The table struct mirrors the database
exactly, including columns the page must never see. Your \code{portfolio\_items} table has a
\code{public} boolean; the DTO for a public listing should not have that field at all, because a
field that does not exist cannot be leaked by a careless \code{view!}.

\textbf{The payload is a published API.} A server function is a public HTTP endpoint. Anyone can
call it. The DTO is the contract you are publishing, and it deserves to be designed rather than
inherited from whatever your table happens to look like.
\end{why}

\subsection{What it looks like}

\begin{lstlisting}[style=rs,caption={\file{src/publications/dto.rs} --- ungated, no Diesel anywhere.}]
use serde::{Deserialize, Serialize};

#[derive(Debug, Clone, PartialEq, Serialize, Deserialize)]
pub struct PublicationDto {
    pub id: i32,
    pub title: String,
    pub abstract_text: String,   // renamed: `abs` is a SQL function name and an unclear field name
    pub year: i32,
    pub journal: String,
    pub link: String,
}
\end{lstlisting}

\begin{lstlisting}[style=rs,caption={\file{src/publications/model.rs} --- server only, unchanged in spirit from today.}]
use diesel::prelude::*;

#[derive(Debug, Queryable, Selectable)]
#[diesel(table_name = crate::schema::publication_items)]
#[diesel(check_for_backend(diesel::pg::Pg))]
pub struct PublicationItem {
    pub id: i32,
    pub title: String,
    pub abstract_text: String,
    pub year: i32,
    pub journal: String,
    pub link: String,
}

impl From<PublicationItem> for crate::publications::dto::PublicationDto {
    fn from(row: PublicationItem) -> Self {
        Self {
            id: row.id,
            title: row.title,
            abstract_text: row.abstract_text,
            year: row.year,
            journal: row.journal,
            link: row.link,
        }
    }
}
\end{lstlisting}

Note that \code{Serialize} and \code{Deserialize} have left the model entirely. The model has no
business being serialised --- it never crosses a network. That deletion is not cosmetic; it is what
makes the boundary real.

\begin{pattern}[Put the \code{From} impl next to the model]
The conversion is a server-side concern and depends on the server-side type, so it belongs in
\file{model.rs}, not \file{dto.rs}. Keeping it there means \file{dto.rs} imports nothing from your
crate at all, which is exactly the property that guarantees it is safe for the browser.
\end{pattern}

\section{The mistake that looks like the fix}

There is an obvious-seeming way to avoid writing two structs: keep one struct and gate the derives.

\begin{antipattern}[Gating serde behind \code{cfg\_attr}]
\begin{lstlisting}[style=rs]
#[derive(Debug, Clone)]
#[cfg_attr(feature = "ssr", derive(serde::Serialize, serde::Deserialize))]
#[cfg_attr(feature = "ssr", derive(diesel::Queryable, diesel::Selectable))]
pub struct PublicationItem { /* ... */ }
\end{lstlisting}
This is worse than the original problem, because it fails \emph{later} and more obscurely. The
type now exists in the browser build but without its serde impls, so the failure moves from the
type definition to every use site.
\end{antipattern}

The reported errors are worth showing, because recognising them saves an hour. Research for this
document reproduced them against \code{leptos 0.8.20}: you get four \code{error[E0277]}s per DTO.
Two come from the \code{\#[server]} macro:

\begin{lstlisting}[style=out]
error[E0277]: the trait bound `Dto: serde::Serialize` is not satisfied
  = note: required for `JsonEncoding` to implement `Encodes<Vec<Dto>>`
  = note: required for `Vec<Dto>` to implement
          `IntoRes<Post<JsonEncoding>, BrowserMockRes, ServerFnError>`
  = note: required by a bound in `leptos::server_fn::ServerFn::Protocol`
\end{lstlisting}

and two more come from \code{Resource::new} independently of the server function, anchored at
\code{leptos\_server-0.8.7/src/resource.rs:1030}, where the bound is
\code{JsonSerdeCodec: Encoder<T> + Decoder<T>}.

\begin{pitfall}[Read the type names in the error]
\code{BrowserMockRes}, \code{BrowserClient}, \code{BrowserMockServer}. Those names are the
diagnostic. They tell you the bound is being checked \emph{in the client build}, which tells you
the missing impl has to exist in the client build, which tells you the \code{cfg\_attr} is the
bug. Whenever a confusing trait error mentions a \code{Browser*} type, ask what is missing on the
client side.
\end{pitfall}

Only \emph{server-only} derives belong in \code{cfg\_attr}. If you ever do collapse a model and a
DTO into one type --- reasonable for a slice that is a pure pass-through --- it looks like this,
with serde unconditional:

\begin{lstlisting}[style=rs]
#[derive(Debug, Clone, serde::Serialize, serde::Deserialize)]        // always
#[cfg_attr(feature = "ssr", derive(diesel::Queryable, diesel::Selectable))]
#[cfg_attr(feature = "ssr", diesel(table_name = crate::schema::publication_items))]
#[cfg_attr(feature = "ssr", diesel(check_for_backend(diesel::pg::Pg)))]
pub struct Publication { /* ... */ }
\end{lstlisting}

\section{Private fields, and why they are a problem now}

Two of your four models declare their fields private:

\begin{lstlisting}[style=rs,caption={\file{src/job\_experience/model.rs} --- no \code{pub} on any field.}]
pub struct JobExperienceItem {
    id: i32,
    date_from: NaiveDate,
    date_to: Option<NaiveDate>,
    job_title: String,
    accomplishments: String,
    responsabilities: String,
}
\end{lstlisting}

\code{PersonalInformation}, \code{ContactInformation} and \code{JobInstitution} are the same;
\code{PortfolioItem}, \code{Tag} and \code{PublicationItem} correctly use \code{pub}. This
compiles today only because nothing outside those modules reads the fields. The moment a
\code{From} impl in another module, or a page in \file{ui/}, tries to read \code{job\_title}, you
get \code{error[E0616]: field `job\_title` of struct `JobExperienceItem` is private}. This is
defect \dnum{7}.

Two acceptable resolutions, and the choice is a real design decision rather than a formality:

\begin{itemize}
  \item \textbf{Make the fields \code{pub}.} Correct for a row struct. A Diesel model is a
    transparent record of what the database returned; it has no invariants to protect, and adding
    getters for all six fields would be ceremony with no content. Do this.
  \item \textbf{Keep them private and add a constructor plus accessors.} Correct only if the type
    is a real domain aggregate that enforces invariants --- ``\code{date\_to} is either \code{None}
    or on/after \code{date\_from}''. If you want that invariant enforced in Rust rather than only
    by a database \code{CHECK} constraint, then the invariant-holding type is a \emph{third} type,
    distinct from both the row and the DTO, and it deserves the privacy. That is a reasonable thing
    to want later. It is not what these structs are today.
\end{itemize}

Whichever you choose, choose it consistently. Right now the codebase does both, and the difference
is not deliberate --- it is which file was written on which evening. Consistency here is worth
more than either option.

% ===========================================================================
\chapter{Repositories, blocking, and the pool}
\label{ch:blocking}

\section{Why \code{Repository<T>} should not survive}

The current abstraction is four lines:

\begin{lstlisting}[style=rs,caption={\file{src/core/mod.rs}, in full.}]
pub trait Repository<T> {
    fn get_all(&mut self) -> anyhow::Result<Vec<T>>;
    fn get_by_id(&mut self, id: i32) -> anyhow::Result<T>;
}
\end{lstlisting}

It is a familiar shape from other languages, and in Rust it has four problems. Taken one at a time,
because only the last one is fatal:

\subsection*{1. The generic parameter buys nothing}

You cannot write \code{impl<T> Repository<T> for PgStore} and specialise it per entity ---
coherence forbids overlapping impls, and specialisation is unstable. So every implementation is
hand-written anyway, exactly as it is now: four structs, four impls, eight method bodies. The
generic parameter adds a type parameter to every signature and returns nothing. A trait per subject
with the methods that subject actually needs --- \code{PublicationRepository} with
\code{list()}, \code{by\_id()}, \code{by\_year()} --- says more and costs less.

\subsection*{2. Two methods is not an interface, it is a coincidence}

\code{get\_all} and \code{get\_by\_id} happen to be what all four subjects need \emph{today},
because you have not written any interesting queries yet. The first time you want
``publications grouped by year, newest first'' or ``portfolio items with their tags'', it will not
fit, and it will go somewhere else --- an inherent method on the struct, next to the trait impl.
Then the trait describes half of each repository and you have to read both to know what a
repository can do.

\subsection*{3. \code{\&mut self} is a misattribution, and it is contagious}

This one is worth being precise about, because the reasoning that produces it is so plausible.
Diesel query methods take \code{\&mut conn}. So a repository method needs a mutable connection. So
--- it seems to follow --- the repository must be \code{\&mut self}. It does not follow. The
\code{\&mut} belongs to the \emph{connection}, which you obtain locally:

\begin{lstlisting}[style=rs]
pub fn get(&self) -> Result<PooledConnection<M>, Error>   // r2d2::Pool::get -- note &self
\end{lstlisting}

\code{r2d2::Pool} is \code{pub struct Pool<M>(Arc<SharedPool<M>>)}. It is \code{Send},
\code{Sync}, and \code{Clone}, and its \code{Clone} is an atomic refcount bump. Concurrent
immutable access is exactly what it is built for. Declaring the repository \code{\&mut self}
throws that away: every caller now needs exclusive access, which in an async server means putting
the repository behind a \code{Mutex} and serialising every request through it --- destroying the
concurrency the pool exists to provide. This is defect \dnum{8}.

\begin{pattern}[Take \code{\&self}, clone the pool]
\begin{lstlisting}[style=rs]
#[derive(Clone)]                                  // cheap: one Arc bump
pub struct PublicationRepository { pool: DbPool }

impl PublicationRepository {
    pub fn new(pool: DbPool) -> Self { Self { pool } }

    pub async fn list(&self) -> Result<Vec<PublicationItem>, RepoError> {
        let pool = self.pool.clone();             // clone into the 'static closure
        // ... see the next section for the closure
    }
}
\end{lstlisting}
And do not wrap it in \code{Arc<Pool<...>>}. It is already an \code{Arc}; a second layer is pure
indirection.
\end{pattern}

\subsection*{4. The signatures are synchronous, and that is the fatal one}

\code{fn get\_all(\&mut self) -> anyhow::Result<Vec<T>>} is a blocking function. Called from a
server function, it blocks an Actix worker thread. The rest of this chapter is about why that
matters more than it sounds like it should.

\section{Blocking: the bug that no test will find}

\begin{keyidea}
Diesel is synchronous. A server function is asynchronous. Calling the former directly from the
latter does not fail --- it silently converts your concurrent server into a sequential one.
\end{keyidea}

\subsection{The mechanism}

An async runtime multiplexes many tasks onto few threads by requiring that every task yield when
it waits. A blocking call does not yield: it holds its thread for the whole duration of the wait,
so the scheduler on that thread cannot poll any other task. Tokio's own documentation is blunt
about it --- a blocking operation in a task ``would block the entire thread, preventing other tasks
from running''.

Actix makes the consequence sharper than generic Tokio, because each HTTP worker processes its
accepted connections on one thread, and the default worker count is
\code{std::thread::available\_parallelism()}. So with $W$ workers and a query taking $T$ seconds,
exactly $W$ queries progress at a time, and throughput caps near $W/T$ requests per second
regardless of how idle Postgres is.

\begin{diagram}{Four Actix workers, each handling one request. Above: Diesel called directly ---
each worker is held for the whole query and cannot serve anyone else. Below: the query is handed to
the blocking pool --- workers return to the event loop immediately.}{fig:blocking}
\begin{tikzpicture}[x=1mm,y=1mm]
  % ---- BAD ----
  \node[anchor=west,font=\sffamily\scriptsize\bfseries,text=crimson] at (0,6) {direct call --- worker blocked};
  \foreach \i in {0,1,2,3} {
    \node[anchor=east,tag] at (-1,{-\i*7}) {worker \i};
    \draw[draw=rule,line width=0.4pt] (0,{-\i*7-3}) rectangle (120,{-\i*7+3});
    \fill[pattern=north east lines,pattern color=crimson!55,draw=crimson!55,line width=0.4pt]
      (4,{-\i*7-3}) rectangle (96,{-\i*7+3});
  }
  \node[lbl,anchor=west] at (30,0) {blocked waiting on Postgres};
  \node[anchor=west,tag,text=crimson] at (122,{-1.5*7}) {queue grows};

  % ---- GOOD ----
  \node[anchor=west,font=\sffamily\scriptsize\bfseries,text=moss] at (0,-32) {web::block --- worker free};
  \foreach \i in {0,1,2,3} {
    \node[anchor=east,tag] at (-1,{-38-\i*7}) {worker \i};
    \draw[draw=rule,line width=0.4pt] (0,{-38-\i*7-3}) rectangle (120,{-38-\i*7+3});
    \fill[fill=moss!22,draw=moss!55,line width=0.4pt] (4,{-38-\i*7-3}) rectangle (12,{-38-\i*7+3});
    \fill[fill=moss!22,draw=moss!55,line width=0.4pt] (92,{-38-\i*7-3}) rectangle (100,{-38-\i*7+3});
    \draw[flow,draw=slate!60,line width=0.5pt] (12,{-38-\i*7}) -- (92,{-38-\i*7});
  }
  \node[lbl,anchor=west] at (26,-38) {worker polls other tasks while it waits};
  \foreach \i in {0,1,2,3} {
    \node[anchor=east,tag,text=slate] at (-1,{-70-\i*5}) {blk. thread \i};
    \draw[draw=rule,line width=0.4pt] (0,{-70-\i*5-2}) rectangle (120,{-70-\i*5+2});
    \fill[fill=slate!18,draw=slate!45,line width=0.4pt] (12,{-70-\i*5-2}) rectangle (92,{-70-\i*5+2});
  }
  \node[lbl,anchor=west] at (30,-70) {the query actually runs here};
\end{tikzpicture}
\end{diagram}

\subsection{What it looks like when it happens}

Not a crash and not a panic --- which is why it is easy to ship. The symptoms are:

\begin{itemize}
  \item flat, low CPU while p99 latency climbs;
  \item request times that cluster in multiples of your query time;
  \item \emph{static assets and health checks timing out too}, because they share the same
    workers;
  \item no improvement when you raise the r2d2 pool size --- the bottleneck is the worker thread,
    not the pool. This symptom is the giveaway.
\end{itemize}

\section{The fix}

\begin{pattern}[Wrap every query in \code{web::block}]
\begin{lstlisting}[style=rs]
pub async fn list(&self) -> Result<Vec<PublicationItem>, RepoError> {
    use crate::schema::publication_items::dsl::*;
    let pool = self.pool.clone();

    let rows = actix_web::web::block(move || {
        let mut conn = pool.get()?;                   // NOTE: get() is INSIDE the closure
        publication_items
            .select(PublicationItem::as_select())
            .load(&mut conn)
    })
    .await?                                            // BlockingError
    .map_err(RepoError::from)?;                        // the inner error

    Ok(rows)
}
\end{lstlisting}
\end{pattern}

Four details in those twelve lines, each of which is a bug if you get it wrong.

\begin{description}
  \item[\code{pool.get()} goes inside the closure.] \code{get()} is itself a blocking call: it
    parks the calling thread for up to \code{connection\_timeout}, which defaults to
    \emph{30 seconds}. Acquiring the connection before the closure reintroduces the exact bug you
    are fixing, in its worst form. The official Actix Diesel example puts \code{get()} inside the
    closure for precisely this reason.

  \item[There are two \code{Result}s, so there are two \code{?}.]
    \code{web::block(...).await} gives you
    \code{Result<Result<T, E>, BlockingError>}. One \code{?} will not compile; an \code{.unwrap()}
    on the outer one swallows the real database error.

  \item[Never hold a \code{PooledConnection} across an \code{.await}.] Acquire it and drop it
    inside one closure. Holding it pins a pooled connection for the whole async section, and
    connections are the scarcest resource you have.

  \item[Blocking tasks are not cancellable.] If the visitor closes the tab, Actix drops the
    request future --- but the closure already handed to the blocking pool runs to completion and
    keeps holding its connection. This is how a pool gets exhausted by clients that already gave
    up, and it is the reason to set a short \code{connection\_timeout}.
\end{description}

\subsection{The alternatives, and why not to take them}

\begin{description}
  \item[\code{tokio::task::spawn\_blocking}] is what \code{web::block} calls internally --- its
    body is literally \code{let fut = actix\_rt::task::spawn\_blocking(f); async \{ fut.await.map\_err(|\_| BlockingError) \}}.
    Note what that \code{map\_err} discards: \code{BlockingError} is a fieldless unit struct, so
    whether the closure panicked or was cancelled is lost. If you want that detail for logging,
    call \code{spawn\_blocking} yourself and keep the \code{JoinError}, which has
    \code{is\_panic()} and \code{into\_panic()}. A reasonable choice for one central helper; not
    worth it at every call site.

  \item[\code{deadpool-diesel}] gives you an async \code{pool.get()} and an \code{interact()}
    method that does the thread hop for you. It adds a dependency to solve a problem
    \code{web::block} already solves, and its newest release (0.6.1) is from May 2024. Skip it
    unless contention on connection \emph{acquisition} becomes your bottleneck, which for a
    personal site it will not.

  \item[\code{diesel-async}] removes the thread hop entirely and is genuinely production-viable
    with Postgres. It is also still 0.x (0.9.2 as of mid-2026), reimplements the Postgres wire
    protocol rather than using \code{libpq}, and has had subtle lifetime bugs --- notably a
    ``Cannot access shared transaction state'' panic when a pooled connection returns to the pool
    before the futures borrowing it are awaited. Adopting it would trade a solved problem for API
    churn with no measurable benefit below thousands of requests per second.
\end{description}

\begin{pattern}[The decision]
Keep synchronous Diesel with r2d2, and wrap every query in \code{web::block}. Revisit only if you
ever measure connection-acquisition contention.
\end{pattern}

\section{Pool sizing}

The defaults will work for a personal site, but two of them are worth changing and the reasoning
generalises.

\begin{lstlisting}[style=rs,caption={A pool built from configuration, with deliberate limits.}]
Pool::builder()
    .max_size(cfg.db_max_connections)             // default is 10
    .connection_timeout(Duration::from_secs(5))   // default is 30 -- far too long
    .test_on_check_out(true)
    .build(ConnectionManager::<PgConnection>::new(&cfg.database_url))
\end{lstlisting}

\code{connection\_timeout} matters because of the non-cancellability point above: a thread waiting
30 seconds for a connection is a blocking-pool slot doing nothing for 30 seconds. Five seconds
fails fast instead.

\code{max\_size} sits inside two independent bounds, and it is easy to reason about only one:

\begin{itemize}
  \item \textbf{Below the blocking capacity.} Actix's blocking pool is \emph{per worker}, not
    global --- \code{worker\_max\_blocking\_threads} defaults to 512 divided by available
    parallelism, so on an eight-core machine that is 64 per worker, 512 in total. Sizing math
    borrowed from plain-Tokio advice (one global 512-thread pool) is wrong under Actix.
  \item \textbf{Below Postgres's budget.} \code{max\_size} multiplied by the number of running
    instances must stay under the server's \code{max\_connections}. r2d2's default of 10 per
    instance quietly exhausts a small managed Postgres as soon as you run two replicas.
\end{itemize}

% ===========================================================================
\chapter{Errors that cross the wire}
\label{ch:errors}

\section{Four layers, four error types}

Error handling is where half-built applications accumulate the most incidental damage, because the
easy thing --- one error type everywhere --- works fine until the day it leaks a connection string
to a browser. The layering below is small enough to adopt now.

\begin{table}[htbp]
\centering
\caption{Which error type belongs where.}
\label{tab:errors}
\small
\begin{tabular}{@{}L{0.18\textwidth}L{0.28\textwidth}L{0.46\textwidth}@{}}
\toprule
\textbf{Layer} & \textbf{Type} & \textbf{Rule} \\
\midrule
driver & \code{diesel::result::Error}, \code{r2d2::Error} & never escapes \file{repository.rs} \\
repository / domain & \code{thiserror} enums, one per subject or use case & named variants for what callers can act on, plus \code{Unknown(anyhow::Error)} \\
wire & one \code{AppError} in \file{src/error.rs} & serde-serialisable, ungated, \emph{carries no internal detail} \\
process & \code{anyhow} & \file{main.rs} and adapter constructors only \\
\bottomrule
\end{tabular}
\end{table}

\begin{why}[Why \code{anyhow} does not belong in the repository]
Your repositories currently return \code{anyhow::Result<T>} with a \code{.context("Can't get
publication items from db")}. That is fine for logging and useless for deciding. A caller cannot
ask an \code{anyhow::Error} ``was this a missing row or a dead connection?'' without matching on
strings. \code{anyhow} is right where the only reasonable response is to log and give up --- which
is \file{main.rs} --- and wrong wherever a caller might branch.

There is also a copy-paste artefact worth noticing: three of the four repositories say
\code{"Can't get publication items from db"} regardless of which table they queried. That is
defect \dnum{11}, and it is a symptom rather than a cause. A typed error would have made the
mistake impossible, because \code{PortfolioError::NotFound} cannot say ``publication''.
\end{why}

\section{The wire error type}

Leptos 0.8 changed the recommended approach here, and tutorials written for 0.6 or 0.7 will lead
you astray. \code{ServerFnError<MyError>} --- the \code{WrappedServerError} variant --- is
\textbf{deprecated as of 0.8.0}. The supported way to return a domain-shaped error is to implement
\code{FromServerFnError} on your own enum.

\begin{lstlisting}[style=rs,caption={\file{src/error.rs} --- ungated, and the only error the browser sees.}]
use leptos::prelude::ServerFnErrorErr;
use leptos::server_fn::codec::JsonEncoding;
use leptos::server_fn::error::FromServerFnError;
use serde::{Deserialize, Serialize};

#[derive(Debug, Clone, Serialize, Deserialize, thiserror::Error)]
pub enum AppError {
    #[error("not found")]
    NotFound,
    #[error("invalid input: {0}")]
    Validation(String),
    #[error("internal error")]
    Internal,                             // deliberately carries nothing
    #[error("transport: {0:?}")]
    ServerFn(ServerFnErrorErr),           // required by the trait below
}

impl FromServerFnError for AppError {
    type Encoder = JsonEncoding;
    fn from_server_fn_error(value: ServerFnErrorErr) -> Self {
        AppError::ServerFn(value)
    }
}
\end{lstlisting}

\begin{pattern}[\code{Internal} must be empty]
\code{AppError::Internal} carries no payload on purpose. Everything unexpected collapses into it,
and the real cause is logged server-side with \code{tracing::error!}. The alternative ---
\code{ServerFnError::ServerError(e.to\_string())} --- sends the string to the browser verbatim,
and \code{e.to\_string()} on a Diesel or r2d2 error can contain your connection string. This is
the single most common way a Rust web application leaks infrastructure detail.
\end{pattern}

\begin{lstlisting}[style=rs,caption={The mapping layer: where a domain error becomes a wire error.}]
#[cfg(feature = "ssr")]
impl From<RepoError> for AppError {
    fn from(e: RepoError) -> Self {
        use actix_web::http::StatusCode;
        let status = expect_context::<leptos_actix::ResponseOptions>();
        match e {
            RepoError::NotFound => {
                status.set_status(StatusCode::NOT_FOUND);
                AppError::NotFound
            }
            other => {
                tracing::error!(error = ?other, "unhandled repository error");
                status.set_status(StatusCode::INTERNAL_SERVER_ERROR);
                AppError::Internal          // nothing leaks
            }
        }
    }
}
\end{lstlisting}

\begin{pitfall}[Server-function errors do not set a status code by themselves]
An error returned from a \code{\#[server]} function travels as an encoded body plus a marker
header (\code{SERVER\_FN\_ERROR\_HEADER}). The HTTP status stays 200 unless you set it. If you
want a genuine 404 --- and you do, for crawlers, caches and uptime monitoring --- you must call
\code{ResponseOptions::set\_status} explicitly, as above.
\end{pitfall}

\begin{pitfall}[\code{?} stops working when you leave the default error type]
With \code{Result<T, ServerFnError>} you get \code{impl<E: std::error::Error> From<E> for
ServerFnError} for free, so \code{?} converts anything. The moment you switch to your own
\code{AppError}, that blanket impl is gone and you must write the conversions. In particular
\code{leptos\_actix::extract()} returns \code{ServerFnErrorErr}, not \code{ServerFnError}, so it
becomes \code{.map\_err(AppError::from\_server\_fn\_error)?}.
\end{pitfall}

\begin{pattern}[Start with \code{ServerFnError}, migrate deliberately]
For the first few server functions, \code{Result<T, ServerFnError>} is fine and keeps \code{?}
ergonomic while you are still finding the shape of the code. Move to \code{AppError} when the UI
first needs to \emph{distinguish} two failures --- typically the first time you want a real 404
page. Do not do both at once; a codebase with two wire-error conventions is worse than either.
\end{pattern}

% ===========================================================================
\chapter{The composition root}
\label{ch:composition}

\section{What is wrong with \file{main.rs} today}

\file{src/database/connection.rs} builds a pool. Nothing calls it. \file{src/main.rs} is the
starter template verbatim: it configures Leptos, serves static files, and never mentions a
database. Closing that gap is the single highest-value change in this document, because nothing
else can be tested until data can reach a page.

There are also three smaller issues in the current file, all in the same few lines:

\begin{lstlisting}[style=rs,caption={\file{src/main.rs}, lines 15--21 --- three problems in seven lines.}]
HttpServer::new(move || {
    let routes = generate_route_list(App);          // (a) fine, but see below
    let leptos_options = &conf.leptos_options;
    let site_root = leptos_options.site_root.clone().to_string();

    println!("listening on http://{}", &addr);      // (b) printed once PER WORKER
    // ...
        .service(Files::new("/assets", &site_root)) // (c) serves the whole site root
\end{lstlisting}

\begin{description}
  \item[(a)] \code{generate\_route\_list(App)} inside the factory is correct --- it is required
    per worker, and it is what the official template does. Mentioned only so you do not ``fix'' it.
  \item[(b)] The \code{println!} is inside the app factory, so it runs once per worker thread and
    you get one line per CPU core at startup. Harmless, confusing, and a sign the line is in the
    wrong scope. Move it above \code{HttpServer::new}. Defect \dnum{17}.
  \item[(c)] \code{Files::new("/assets", \&site\_root)} maps the URL \code{/assets} onto
    \emph{the whole site root}, not onto the assets subdirectory --- so
    \code{/assets/pkg/egonik-site.wasm} is also reachable. It comes from the upstream template and
    is not a security hole (everything under \code{site\_root} is public anyway), but it is
    surprising. Defect \dnum{18}.
\end{description}

\section{\code{AppState}, and the one line that must move}
\label{sec:appstate}

This section was going to propose an \code{AppState}. While the document was being written you
wrote one, so instead it endorses what you wrote and flags a single line that will exhaust your
database.

\begin{verified}[What is in the tree now]
\begin{lstlisting}[style=rs,caption={\file{src/entrypoint/application\_state.rs}, as written.}]
#[derive(Debug, Clone)]
pub struct AppState {
    pub job_experience_repository: JobExperienceItemRepository,
    pub personal_information_repository: PersonalInformationRepository,
    pub portfolio_items_repository: PortfolioItemsRepository,
    pub publications_repository: PublicationsRepository,
}

impl AppState {
    pub fn new(pool: DbPool) -> Self {
        // one pool, cloned into four repositories -- each clone is an Arc bump
        Self {
            personal_information_repository: PersonalInformationRepository::new(pool.clone()),
            job_experience_repository: JobExperienceItemRepository::new(pool.clone()),
            portfolio_items_repository: PortfolioItemsRepository::new(pool.clone()),
            publications_repository: PublicationsRepository::new(pool),
        }
    }
}
\end{lstlisting}
\end{verified}

\begin{pattern}[This is the right shape --- keep it]
Holding the four \emph{repositories} rather than the raw pool is better than what this chapter was
going to suggest. A server function then asks for the capability it needs
(\code{state.publications\_repository}) rather than for a database handle it has to know what to do
with, and the pool stops being ambient. Two details you got right and that are worth knowing you
got right:

\code{Clone} is required, because Actix hands application data to each worker, and it is cheap
because every repository is one \code{Arc} bump --- which is why \code{\#[derive(Clone)]} on
\code{PublicationsRepository} was the correct accompanying change.

Taking \code{pool} as a parameter rather than calling \code{get\_connection\_pool()} inside
\code{new} keeps the dependency in the signature, which is exactly the property
\S\ref{sec:config} argues for.
\end{pattern}

Two things to add when convenient: \code{config: Arc<AppConfig>} alongside the repositories, so that
things like a page size or a base URL stop being constants; and \code{new} returning
\code{anyhow::Result<Self>} once the pool construction stops panicking.

\subsection{The line that must move}

\begin{antipattern}[The pool is being built inside the worker factory]
\begin{lstlisting}[style=rs,caption={\file{src/main.rs}, as written. The closure body runs
\emph{once per worker}.}]
HttpServer::new(move || {
    let pool = get_connection_pool();          // <-- HERE. Once per worker.
    let routes = generate_route_list(App);
    let app_state = AppState::new(pool);
    println!("listening on http://{}", &addr);

    App::new()
        .app_data(web::Data::new(app_state))
        // ...
})
\end{lstlisting}
\end{antipattern}

Everything in that closure is constructed once per worker thread, and you already have proof of how
many times that is: the \code{println!} on the following line printed \textbf{twelve times} when the
binary was run for this document. Twelve invocations of the factory means twelve independent
connection pools.

\begin{verified}[Verified: this over-subscribes Postgres by 20\%]
\begin{lstlisting}[style=out]
$ nproc                                     # = Actix's default worker count
12
$ grep -c 'listening on' server.log         # = times the factory closure ran
12

r2d2's Builder::max_size default              = 10   (connection.rs sets no max_size)
pools created                                 = 12
maximum connections this process will open     = 120

$ psql -tAc 'show max_connections'
100
\end{lstlisting}
\end{verified}

So under load this process alone tries to open 120 connections against a server that permits 100,
and the failures arrive as \code{r2d2} timeouts surfacing as 500s --- after a 30-second wait each,
because \code{connection\_timeout} is also at its default. Note that it will not fail in
development: each pool opens connections lazily, and one browser tab never needs more than a handful.
That is what makes it worth fixing before it can be observed.

\begin{pattern}[Build it once, above the factory, and clone in]
\begin{lstlisting}[style=rs]
let pool = get_connection_pool();          // once
let app_state = AppState::new(pool);       // once
println!/tracing::info!(...);              // once

HttpServer::new(move || {
    let routes = generate_route_list(App); // this one genuinely IS per-worker
    App::new()
        .app_data(web::Data::new(app_state.clone()))   // one Arc bump per worker
        // ...
})
\end{lstlisting}
\end{pattern}

\begin{why}[Why \code{generate\_route\_list} stays inside and the pool does not]
They look symmetrical and they are not. \code{generate\_route\_list(App)} produces a
\code{Vec<ActixRouteListing>} that \code{leptos\_routes} consumes when building \emph{that worker's}
service tree, so it belongs to the factory --- and it allocates nothing scarce. The pool holds
operating-system sockets against a server with a hard limit. The rule is not ``move everything
out of the factory''; it is \textbf{anything holding a finite external resource is constructed
once}.
\end{why}

This is defect \dnum{38}, and it is the only new blocker introduced since the register was written.

\section{Getting the pool into a server function}

Two mechanisms exist. Both work; they differ in what they couple you to.

\begin{diagram}{Two routes from \code{main} to a server-function body. The Actix route (left) is
what the Leptos book documents; the context route (right) is what you want once a server function
should receive a domain service rather than an Actix type.}{fig:wiring}
\begin{tikzpicture}[node distance=6mm]
  \node[srv,minimum width=44mm] (main) {\file{main.rs}\\[1pt]\scriptsize builds the pool \emph{once}};

  \node[srv,minimum width=50mm,below left=10mm and 2mm of main] (ad)
    {\code{.app\_data(web::Data::new(pool.clone()))}};
  \node[srv,minimum width=50mm,below right=10mm and 2mm of main] (ctx)
    {\code{.leptos\_routes\_with\_context(routes,}\\\code{\ \ move || provide\_context(pool.clone()), ...)}};

  \node[shared,minimum width=50mm,below=9mm of ad] (ex)
    {\code{leptos\_actix::extract::<Data<Pool>>()}\\[1pt]\scriptsize inside the \code{\#[server]} body};
  \node[shared,minimum width=50mm,below=9mm of ctx] (uc)
    {\code{use\_context::<Pool>()}\\[1pt]\scriptsize inside the \code{\#[server]} body};

  \node[box,fill=wash,minimum width=104mm,below=9mm of $(ex)!0.5!(uc)$] (body)
    {the query runs, wrapped in \code{web::block}};

  \draw[flow] (main) -- (ad);
  \draw[flow] (main) -- (ctx);
  \draw[flow] (ad) -- (ex);
  \draw[flow] (ctx) -- (uc);
  \draw[flow] (ex) -- (body);
  \draw[flow] (uc) -- (body);
\end{tikzpicture}
\end{diagram}

\subsection{The Actix route --- start here}

\begin{lstlisting}[style=rs,caption={The wiring, in \file{main.rs}.}]
let pool = build_pool(&cfg)?;                     // built ONCE, outside the factory

HttpServer::new(move || {
    App::new()
        .app_data(web::Data::new(pool.clone()))   // one Arc bump per worker
        .app_data(web::Data::new(leptos_options.to_owned()))
        .service(Files::new("/pkg", format!("{site_root}/pkg")))
        .service(favicon)
        .leptos_routes(routes, { /* the HTML shell */ })
})
\end{lstlisting}

\begin{lstlisting}[style=rs,caption={And reading it back, in \file{src/publications/api.rs}.}]
#[server]
pub async fn list_publications() -> Result<Vec<PublicationDto>, ServerFnError> {
    use actix_web::web::Data;
    use crate::database::connection::DbPool;
    use crate::publications::repository::PublicationRepository;

    let pool: Data<DbPool> = leptos_actix::extract().await?;
    let repo = PublicationRepository::new(pool.get_ref().clone());
    Ok(repo.list().await?.into_iter().map(Into::into).collect())
}
\end{lstlisting}

\begin{why}[Why no extra route registration is needed]
A reasonable expectation, carried over from earlier Leptos versions, is that server functions need
their own Actix route --- something like
\code{.route("/api/\{tail:.*\}", leptos\_actix::handle\_server\_fns())}. In
\code{leptos\_actix} 0.8.7 that is \textbf{not} required: \code{leptos\_routes} registers every
path from \code{server\_fn::actix::server\_fn\_paths()} itself, deliberately \emph{before} the
router's wildcard route --- the source carries the comment ``register server functions first to
allow for wildcard route in Leptos's Router''. \code{leptos\_routes} is literally
\code{leptos\_routes\_with\_context(paths, || \{\}, app\_fn)}.

Adding a second mount is not just redundant: if the duplicate lacks the same context closure it
will silently drop context for browser-initiated calls, which produces a bug that appears only
after hydration and never during SSR. Do not add it.
\end{why}

\begin{pitfall}[\code{extract()} sees the request head, not the body]
\code{leptos\_actix::extract} is implemented as \code{use\_context::<Request>()} followed by
\code{T::extract(\&req)}. The body has already been consumed by the server-function decoder, so
\code{web::Json<T>} and \code{web::Form<T>} extractors will fail. \code{web::Data<T>},
\code{Query<T>}, \code{ConnectionInfo}, \code{HttpRequest} and headers all work.
\end{pitfall}

\subsection{The context route --- where you will end up}

The Actix route works, and it couples your server-function bodies to \code{actix\_web::web::Data}.
That is a small cost now and an annoying one later, if you ever want to call a server function from
a test without an Actix request in scope. The alternative injects your own type:

\begin{lstlisting}[style=rs]
let pool_for_ctx = pool.clone();
App::new()
    .leptos_routes_with_context(
        routes,
        move || provide_context(pool_for_ctx.clone()),   // both SSR *and* server-fn routes
        { /* the HTML shell */ },
    )
\end{lstlisting}

\begin{lstlisting}[style=rs]
let pool = use_context::<DbPool>()
    .ok_or_else(|| ServerFnError::new("pool missing from context"))?;
\end{lstlisting}

\begin{why}[Why this is safe in 0.8.7 and was not always]
\code{leptos\_routes\_with\_context} registers the server-function routes with the \emph{same}
context closure it uses for rendering. That matters because a server function is reached two
different ways --- called directly by the renderer during SSR, and posted to over HTTP after
hydration --- and both paths must have the context. The \code{leptos\_actix} docs warn about
exactly this: if you mount \code{handle\_server\_fns} by hand you must provide the context to both
handlers yourself. Using \code{leptos\_routes\_with\_context} and nothing else avoids the whole
class of mistake.
\end{why}

\begin{pattern}[The order to do it in]
Use \code{app\_data} plus \code{extract} for the first server function, because it is three lines
and gets data on screen today. Switch to \code{leptos\_routes\_with\_context} when you introduce a
service type that owns several repositories --- at that point you are injecting your own
abstraction rather than an Actix wrapper, and the context route reads better.
\end{pattern}

\section{Configuration, and the \code{dotenv} call in the wrong place}
\label{sec:config}

\begin{lstlisting}[style=rs,caption={\file{src/database/connection.rs} as it stands.}]
pub fn get_connection_pool() -> DbPool {
    dotenv().ok();                                              // <-- (1)
    let database_url = env::var("DATABASE_URL").expect("...");  // <-- (2)
    // ...
    .build(manager).expect("Could not build connection pool")    // <-- (3)
}
\end{lstlisting}

Three things to change, and the third is the one that will actually hurt.

\begin{enumerate}
  \item \textbf{\code{dotenv()} inside a constructor.} It mutates process-global state from deep
    inside a data-access function. Call it once, at the top of \code{main}, as the dev-time
    convenience it is. Defect \dnum{12}.

  \item \textbf{Reading the environment inside a constructor} hides a dependency from the type
    signature and makes the function untestable without mutating the process environment --- which
    in Rust 2024 is genuinely awkward, since \code{std::env::set\_var} is \code{unsafe} because it
    is not thread-safe. Pass a value in: \code{fn build\_pool(cfg: \&AppConfig) -> anyhow::Result<DbPool>}.

  \item \textbf{\code{.expect()} in a library function.} A panic here aborts the process with a
    message and no context. Return \code{anyhow::Result} and let \code{main} decide; \code{main}
    is where ``log it and exit non-zero'' is the correct response.
\end{enumerate}

\begin{lstlisting}[style=rs,caption={\file{src/config.rs} --- typed, loaded once, validated at the edge.}]
#[derive(Debug, serde::Deserialize)]
pub struct AppConfig {
    pub database_url: String,
    #[serde(default = "default_pool_size")]
    pub db_max_connections: u32,
}

impl AppConfig {
    pub fn load() -> anyhow::Result<Self> {
        Ok(config::Config::builder()
            .add_source(config::File::with_name("config/base").required(false))
            .add_source(config::Environment::with_prefix("EGONIK").separator("__"))
            .build()?
            .try_deserialize()?)
    }
}
fn default_pool_size() -> u32 { 8 }
\end{lstlisting}

\begin{pitfall}[Do not use the \code{LEPTOS\_} prefix for your own variables]
\code{LEPTOS\_} belongs to \code{leptos\_config}. Every field of \code{LeptosOptions} ---
\code{site\_root}, \code{site\_addr}, \code{reload\_port}, \code{hash\_files}, and a dozen more ---
maps to a \code{LEPTOS\_}-prefixed variable, and a future field could collide with yours. Use your
own prefix. \code{EGONIK\_DATABASE\_URL} in the snippet above.
\end{pitfall}

\begin{pitfall}[\code{get\_configuration(Some("./Cargo.toml"))} reads Cargo.toml at runtime]
Keep the \code{None} that the template gave you. Passing a path makes the running binary read
\file{Cargo.toml} on startup, which works in development and fails in a container that ships only
the binary and the site directory. With \code{None}, \code{LeptosOptions} comes from
\code{LEPTOS\_*} environment variables --- which is exactly what a Dockerfile should set.
\end{pitfall}

\section{The whole composition root}

Putting it together, this is what \file{main.rs} should look like. It is longer than the template
and it is the only place in the codebase that knows how the pieces connect.

\begin{lstlisting}[style=rs,caption={The target \file{src/main.rs}.}]
#[cfg(feature = "ssr")]
#[actix_web::main]
async fn main() -> anyhow::Result<()> {
    use actix_files::Files;
    use actix_web::{web, App, HttpServer};
    use leptos::config::get_configuration;
    use leptos::prelude::*;
    use leptos_actix::{generate_route_list, LeptosRoutes};
    use leptos_meta::MetaTags;
    use egonik_site::{app::*, config::AppConfig, database::connection::build_pool};

    dotenvy::dotenv().ok();                       // once, here, dev convenience only
    tracing_subscriber::fmt::init();              // see Chapter on operations

    let cfg = AppConfig::load()?;                 // typed and validated before anything starts
    let pool = build_pool(&cfg)?;                 // fail fast if the database is unreachable

    let conf = get_configuration(None)?;          // LEPTOS_* only
    let addr = conf.leptos_options.site_addr;
    tracing::info!(%addr, "starting server");     // ONCE, not per worker

    HttpServer::new(move || {
        let routes = generate_route_list(App);
        let leptos_options = &conf.leptos_options;
        let site_root = leptos_options.site_root.clone().to_string();

        App::new()
            .app_data(web::Data::new(pool.clone()))
            .app_data(web::Data::new(leptos_options.to_owned()))
            .service(Files::new("/pkg", format!("{site_root}/pkg")))
            .service(favicon)
            .leptos_routes(routes, {
                let leptos_options = leptos_options.clone();
                move || view! {
                    <!DOCTYPE html>
                    <html lang="en">
                        <head>
                            <meta charset="utf-8"/>
                            <meta name="viewport" content="width=device-width, initial-scale=1"/>
                            <AutoReload options=leptos_options.clone()/>
                            <HydrationScripts options=leptos_options.clone()/>
                            <MetaTags/>
                        </head>
                        <body><App/></body>
                    </html>
                }
            })
    })
    .bind(&addr)?
    .run()
    .await?;
    Ok(())
}
\end{lstlisting}

\begin{pattern}[Fail fast at startup]
\code{build\_pool} returning \code{Result} into a \code{main} that returns
\code{anyhow::Result<()>} means an unreachable database stops the process immediately with a
readable message, instead of producing a 500 on the first request. For a personal site running
under a supervisor that restarts it, this is the behaviour you want: a container that will not
start is far easier to diagnose than one that starts and silently serves errors.
\end{pattern}
